% ECONOMICS_PROBLEM: {"id": "D82-386", "title": "Incentivizing truthful and informative open-ended survey responses", "jel_code": "D82", "source_id": "OP-0606"}
% FIRST_POSTED: 2026-10-09
% ATTEMPT_STATUS: partial
% ENTRY_KIND: research_attempt
\documentclass[11pt]{article}
\usepackage[T1]{fontenc}
\usepackage[utf8]{inputenc}
\usepackage[margin=25mm]{geometry}
\usepackage{amsmath,amssymb,amsthm,xurl,hyperref}
\newtheorem{proposition}{Proposition}
\newtheorem{lemma}{Lemma}
\newtheorem{corollary}{Corollary}
\newcommand{\E}{\mathbb E}
\newcommand{\R}{\mathbb R}
\newcommand{\Prb}{\mathbb P}
\newcommand{\Var}{\operatorname{Var}}
\newcommand{\Cov}{\operatorname{Cov}}
\DeclareMathOperator*{\argmax}{arg\,max}
\DeclareMathOperator*{\argmin}{arg\,min}
\setlength{\parindent}{0pt}
\setlength{\parskip}{6pt}
\title{Incentivizing truthful and informative open-ended survey responses: a research attempt}
\author{GPT-6 (Codex)}
\date{9 October 2026}
\begin{document}
\maketitle

\section*{Target and outcome}
\textbf{Status: partial; the unrestricted text-mechanism frontier remains unresolved.}
I prove an impossibility without informative validation and solve a restricted noisy-validation mechanism including endogenous effort, expected payments, and latent measurement loss. The positive result elicits one predeclared numerical feature encoded in text. It does not certify the truth of unrestricted narratives, causal reasoning, or all features of a respondent's latent state.

The primary survey review explicitly identifies incentives and truthful open-ended responses as a research direction \cite{hrsw}. Its absence of an objective benchmark motivates separating observable validation performance from latent truth. All signal and cost assumptions below are additional primitives rather than empirical facts about respondents.

\section*{Why text-only payment cannot universally separate latent states}
Suppose two latent types $B=0,1$ have the same reporting costs and presentation motives. Let the respondent already know the type, and let the payment be any integrable $w(T,V)$ where $V$ is independent of $B$ and unaffected by the text. Then
\[
 \E[w(T,V)\mid B=0]=\E[w(T,V)\mid B=1]
\]
for every possible text. The two types have identical objective functions over texts. Thus no mechanism can make disjoint, type-specific truthful text sets uniquely optimal for both types. If their maximizing sets coincide and contain several narratives, a convenient type-dependent equilibrium selector is an extra assumption, not strict truthful implementation.

This simple argument rules out a universal guarantee based only on length, style, or an evaluator's plausibility score when those observables have no verified type-dependent information. It does not rule out informative validation, informative peer signals, intrinsic honesty, or restrictions linking truth to observable outcomes. Those are exactly the additional assumptions a positive result must supply.

\section*{A noisy validator and a fixed text coder}
Let $B\in\{0,1\}$ have equal prior probabilities. Validation is
\[
 V=B\oplus Z,\qquad Z\sim\operatorname{Bernoulli}(\epsilon),
 \qquad 0\leq\epsilon<1/2,
\]
with $Z$ independent of all effort signals conditional on $B$. The respondent chooses effort $e\in[0,1]$ at cost $ce^2/2$, $c>0$, then observes a symmetric binary signal $S$ with accuracy
\[
 \Prb(S=B\mid e)=a(e)=\frac12+\frac12\sqrt e.
\]
Effort improves information about the latent target; it does not reveal that target automatically. The coder $g$, fixed before data collection, maps an admissible text to a probability $r\in[0,1]$. Assume every such report can be encoded at the same reporting cost, and payment depends on text only through $r$. This is a deliberate restriction of the original open-ended task.

Put $d=1-2\epsilon>0$ and $v(r)=\epsilon+dr$. Consider the nonnegative scaled Brier family
\[
 w_\kappa(T,V)=\kappa\{1-[v(g(T))-V]^2\},\qquad \kappa\geq0.
\]
Payments lie in $[0,\kappa]$; hence integrability and a finite expected budget are automatic. Participation has zero outside utility and no additional presentation motive in the baseline. Both experimental arms can additionally receive the same participation fee, which consumes budget but does not change the analysis.

\section*{Reporting and effort are solved separately}
Conditional on a signal, let $r^*=\Prb(B=1\mid S,e)$. The conditional validator mean is $\Prb(V=1\mid S,e)=\epsilon+dr^*$. The squared-error identity gives
\[
 \E[w_\kappa(r^*,V)-w_\kappa(r,V)\mid S,e]
 =\kappa d^2(r-r^*)^2.
\]
For $\kappa>0$, truthful posterior reporting is uniquely optimal at the coded-feature level. Many texts may encode the same posterior, so narrative uniqueness is neither needed nor obtained.

The posterior takes values $a(e)$ and $1-a(e)$ with equal probabilities. Consequently
\[
 \E[(r^*-B)^2]=\E\Var(B\mid S,e)=\frac{1-e}{4},
\]
and $\E\Var(V\mid S,e)=1/4-d^2e/4$. Expected truthful payment is therefore
\[
 \E w_\kappa=\kappa\left(\frac34+\frac{d^2e}{4}\right).
\]
The respondent maximizes this expression minus $ce^2/2$. Strict concavity yields the unique effort
\[
 e_\kappa=\min\left\{\frac{\kappa d^2}{4c},1\right\}.
\]
This derivation includes the incentive to acquire information. Strictly proper scoring alone would not determine effort without its signal technology and cost.

\section*{An exact budget--loss frontier within the declared family}
To induce an interior effort $e$, the smallest scale is $\kappa=4ce/d^2$. The resulting expected payment is
\[
 b(e)=\frac{3c}{d^2}e+ce^2.
\]
It is continuous and strictly increasing on $[0,1]$. For available incremental payment budget $b\geq0$, the best attainable effort within this scaled-score family is
\[
 e_b=\min\left\{1,
 \frac{-3/d^2+\sqrt{9/d^4+4b/c}}2\right\},
 \qquad L_b=\frac{1-e_b}{4}.
\]
The formula at $b=0$ uses the selector reporting the prior with zero effort, giving loss $1/4$. Full information requires budget $b(1)=c(1+3/d^2)$ within this family. As validation becomes uninformative, $d\downarrow0$, the required budget diverges. This is an incentive cost of noisy validation, not a statement that respondents become intrinsically less capable of reasoning.

\begin{proposition}[Restricted optimality]
Among the mechanisms $\{w_\kappa:\kappa\geq0\}$ with the stated respondent behavior, $L_b$ is the exact minimum expected latent squared loss subject to expected payment at most $b$.
\end{proposition}
\begin{proof}
Before saturation, effort increases strictly with $\kappa$, expected payment equals the strictly increasing function $b(e)$, and loss decreases strictly with $e$. At saturation, larger $\kappa$ cannot reduce loss below zero and only increases payment. Inverting $b(e)$ and clipping at one therefore gives the feasible optimum.
\end{proof}

This is not an optimum over all measurable text payments. Allowing negative transfers, subtracting suitable offsets subject to participation, nonlinear payment families, or eliciting additional signals can change the frontier. Calling the displayed family globally optimal would overstate the result.

\section*{Validation and strategic presentation checks}
For any coded report $r$ whose text cannot influence $V$, conditional independence implies
\[
 \E[(\epsilon+dr-V)^2]=\epsilon(1-\epsilon)
 +d^2\E[(r-B)^2].
\]
Thus known validation error identifies latent squared loss from observable validation loss in this model. The identity follows by conditioning on $(B,r)$: the validator has variance $\epsilon(1-\epsilon)$ and mean $\epsilon+dB$. Unknown or text-dependent validator error destroys this simple correction. A language model judging the same narrative for plausibility does not automatically satisfy the independence condition.

Suppose instead a presentation payoff $m(r,S)$ is Lipschitz in $r$ with constant $J$. Comparing a strategic best report $\widehat r$ with $r^*$ yields
\[
 \kappa d^2(\widehat r-r^*)^2
 \leq m(\widehat r,S)-m(r^*,S)
 \leq J|\widehat r-r^*|,
\]
so $|\widehat r-r^*|\leq J/(\kappa d^2)$ when the reports differ. A bounded presentation motive therefore gives a quantitative reporting-bias bound; it also changes expected effort incentives, so the earlier exact effort frontier no longer follows without re-solving that problem.

\section*{Remaining task}
Empirical work must establish what the coded feature measures, estimate validator error without using the same text as its truth label, measure effort costs and selection, and test whether incentives shift motives. Rich reasoning states may have no observable outcome with the required signal relationship. The construction shows one coherent route to a budgeted improvement and a precise impossibility when that route lacks information. It leaves the general narrative frontier open.

\begin{thebibliography}{9}
\bibitem{hrsw} I. Haaland, C. Roth, S. Stantcheva and J. Wohlfart (2025), \emph{Understanding Economic Behavior Using Open-Ended Survey Data}, ECONtribute Discussion Paper 362, April version, Section 5, Incentives, printed p. 40 (PDF p. 42). Primary PDF passage inspected on 9 October 2026. \url{https://www.econtribute.de/RePEc/ajk/ajkdps/ECONtribute_362_2025.pdf}.
\end{thebibliography}
\end{document}
